\section{Computations}\label{sec:numerics}

All results of this section are numerical. They are a first three-dimensional implementation, run on a shared machine with 7\,GB of memory and two cores; the finer levels needed for asymptotic rates are listed in Appendix~\ref{app:repro} with the exact commands.

\subsection{Solver and meshes}
The solver (\texttt{code3d/svn3d.py}) uses continuous $P_k$ velocities, $k=3,4$, and discontinuous full-$P_{k-1}$ pressures on every sub-tetrahedron of an Alfeld refinement. The forms are those of Section~\ref{sec:setting}, with $h$ the largest tetrahedron diameter and $n$ pointing into the obstacle on $\Gh$; $\GS$ omits and $\CNS$ includes the mean-free traction. The pressure is eliminated by the iterated-penalty method (penalty $10^4$; the velocity matrix is factorised once with MKL PARDISO, and the $\CNS$ traction enters only the right-hand side), stopped when $\norm{\div u_h}<10^{-13}$ and the momentum residual is below $10^{-11}$, or at round-off stagnation. A nonsymmetric LU variant gives the same $u_h$ to $10^{-12}$.

\emph{Sphere meshes.} An equiangular $N\times N$ grid on each face of the box $R=(-2,2)^3$ is projected radially onto the unit sphere; the projected points are the vertices of a convex inscribed polyhedron $P_h$ with $12N^2$ faces. Radial layers at $r_j=|b|^{j/m}$, $m=N/2$, between $\partial P_h$ and $\partial R$ give prisms, which are split into three tetrahedra by a minimum-index rule and then Alfeld-refined. Every boundary face is a whole face of $P_h$ or lies on $\partial R$. Shape regularity is uniform ($\min\text{vol}/\diam^3=0.0074$--$0.0079$ for $N=2,\dots,24$), but $h/\hG=3.26,4.14,4.79,5.24$ for $N=2,4,6,8$: the family is not self-similar at small $N$, which contaminates rates measured in $\hG$.

\emph{Test problems.} (box) the cube $(-1,1)^3$ with a smooth manufactured solution and fully consistent Nitsche conditions on all faces; (P) $u=0$, $f=\nabla p$ with the wall pressure of the next test, so that $\ut=0$, $\Gc\equiv0$ and the $\GS$ output is exactly the traction response; (sph) Stokes flow past the fixed unit sphere with uniform stream $U=(1,0.4,-0.3)$, with nonconstant wall pressure; (rot) $u=(1-r^{-3})\,\omega\times x$, $p=0$, so that $u=0$ and $p$ is constant on $\Gamma$ ($G=0$).

\subsection{Verification}
\emph{Patch test.} A $P_k$ solution is reproduced to $10^{-9}$--$10^{-11}$ in $H^1$ for $k=3,4$, with Nitsche and with strong conditions. \emph{Orientation.} On the $N=2$ sphere mesh the computed normal points into $P_h$ on $\Gh$ and outward on the box (checked independently).

\begin{center}\small
\begin{tabular}{rrrcccccc}
\toprule
$k$ & $N$ & vel.\ dofs & $h$ & $H^1$ error & rate & $L^2$ error & rate & $\norm{\div u_h}$\\
\midrule
3&2&3\,189&1.732&$3.015\cdot10^{-1}$&--&$2.439\cdot10^{-2}$&--&$1.3\cdot10^{-13}$\\
3&3&10\,290&1.155&$9.780\cdot10^{-2}$&2.78&$5.265\cdot10^{-3}$&3.78&$1.7\cdot10^{-13}$\\
3&4&23\,871&0.866&$4.298\cdot10^{-2}$&2.86&$1.677\cdot10^{-3}$&3.98&$2.3\cdot10^{-13}$\\
3&5&46\,038&0.693&$2.252\cdot10^{-2}$&2.90&$6.812\cdot10^{-4}$&4.04&$3.0\cdot10^{-13}$\\
3&6&78\,897&0.577&$1.322\cdot10^{-2}$&2.92&$3.250\cdot10^{-4}$&4.06&$3.9\cdot10^{-13}$\\
4&2&7\,083&1.732&$3.820\cdot10^{-2}$&--&$3.266\cdot10^{-3}$&--&$6.7\cdot10^{-13}$\\
4&3&23\,115&1.155&$8.207\cdot10^{-3}$&3.79&$4.653\cdot10^{-4}$&4.81&$1.0\cdot10^{-12}$\\
4&4&53\,907&0.866&$2.649\cdot10^{-3}$&3.93&$1.138\cdot10^{-4}$&4.90&$1.5\cdot10^{-12}$\\
\bottomrule
\end{tabular}
\end{center}
\noindent\textbf{Table 1.} Box test, $\mu=400$. The rates approach $k$ in $H^1$ and $k+1$ in $L^2$, and $u_h$ is divergence-free to round-off.

\emph{Stability.} For the pair $(\VO,\Pih\cap L^2_0)$ with the vector Laplacian, a dense eigenvalue computation finds exactly one zero eigenvalue (the constants) and no spurious modes in every case; the discrete inf-sup constant is $0.142$, $0.149$, $0.150$ on the box for $N=1,2,3$, and $0.069$ on the sphere mesh $N=2$. Uniformity in $h$ is therefore seen only on the box.

\emph{Coercivity.} The threshold $\mu_Z^\ast$ of $N_h$ on $\Zh$ (from the inertia of the penalised matrix, independent of the penalty $10^5$--$10^7$, validated against a dense eigenvalue computation) is $76.3$ ($N=2$) and $148.0$ ($N=4$) for $k=3$, and $193.6$ ($N=2$) for $k=4$. It grows with $h/\hG$, as $\mu_0\propto\rho$ predicts (Section~\ref{sec:penalty}). A local sufficient threshold on $\Vh$, computed cell by cell, is $216$, $341$, $421$, $476$ for $N=2,4,6,8$ ($k=3$). All sphere runs below use $\mu\ge200>\mu_Z^\ast$.

\subsection{The leak law}
\begin{center}\small
\begin{tabular}{rrcccccc}
\toprule
$\mu$ & $N$ & $\norm{\nabla(\ut-u_h)}$ & $\dfrac{\norm{\nabla(\ut-u_h)}}{(h/\mu)G}$ & slip & leak & energy & $\norm{\div u_h}$\\
\midrule
200&2&$1.34\cdot10^{-1}$&2.34&0.913&0.891&1.001&$1.0\cdot10^{-15}$\\
200&4&$9.21\cdot10^{-2}$&2.41&0.949&0.943&1.012&$1.4\cdot10^{-15}$\\
400&2&$7.26\cdot10^{-2}$&2.54&0.941&0.923&0.992&$8.7\cdot10^{-16}$\\
400&4&$4.78\cdot10^{-2}$&2.50&0.970&0.965&1.002&$7.5\cdot10^{-16}$\\
1000&2&$3.30\cdot10^{-2}$&2.90&0.963&0.949&0.991&$2.3\cdot10^{-16}$\\
1000&4&$2.03\cdot10^{-2}$&2.66&0.984&0.981&1.001&$3.0\cdot10^{-16}$\\
\bottomrule
\end{tabular}
\end{center}
\noindent\textbf{Table 2.} Test P, $\GS$, $k=3$. slip $:=\norm{\ut-u_h}_{\Gh}/((h/\mu)G)$, leak $:=\ip{u_h\cdot n}{p-\pbar}_{\Gh}/((h/\mu)G^2)$, energy $:=\tnorm{\ut-u_h}(\mu/h)^{1/2}/G$, with $G$ computed on $\Gh$. $N=2$ and $N=4$ have $9\,096$ and $70\,026$ velocity unknowns.

The $\GS$ error is of order $h/\mu$ (the $H^1$ ratio is roughly constant, near $2.5$), and the slip, leak and energy constants approach $1$, with $1-\text{slip}$ decreasing with $h/\mu$: this is Corollary~\ref{er:Bp} in three dimensions. The $H^1$ ratio is not universal (in two dimensions it converges to a domain-dependent constant \cite{PadhiLong2D}). $\CNS$ has no leak: at $\mu=400$ its error on test P is $2.7\cdot10^{-3}$ ($N=2$) and $2.6\cdot10^{-4}$ ($N=4$), at the level of the pressure projection. At $\mu=200$ the $N=4$ $\CNS$ solve returns an error of order $1$: $\mu=200$ is below the well-posedness threshold $\gamma_2$ of Theorem~\ref{er:D} on that mesh, although it exceeds $\mu_Z^\ast$. We report $\CNS$ only for $\mu\ge400$.

\subsection{Sphere tests}
\begin{center}\small
\begin{tabular}{llrcccccc}
\toprule
test & $N$ & $\hG$ & $\GS$ $H^1$ & $\CNS$ $H^1$ & interp.\ $H^1$ & $\GS$ $\norm{e}_{\Gh}$ & $D$: $H^1$ & $D$: slip / leak\\
\midrule
sph&2&0.919&1.269&1.265&1.633&0.445&0.138&1.254 / 0.978\\
sph&4&0.524&0.338&0.325&0.192&0.110&0.064&0.978 / 0.925\\
rot&2&0.919&2.209&2.225&2.390&0.867&--&--\\
rot&4&0.524&0.541&0.549&0.251&0.223&--&--\\
\bottomrule
\end{tabular}
\end{center}
\noindent\textbf{Table 3.} Tests sph and rot, $k=3$, $\mu=400$. ``interp.'' is the $H^1$ error of the nodal interpolant of $\ut$; $D:=u_h^{\GS}-u_h^{\CNS}$, with slip and leak ratios defined as in Table~2. $\norm{\div u_h}$ is $10^{-13}$--$10^{-10}$ in all runs.

\subsection{Interpretation}
Confirmed at these levels: pointwise incompressibility; the optimal rates $h^k$ and $h^{k+1}$ on the box for $k=3,4$; a stable pair without spurious pressure modes; the growth of the coercivity threshold with $h/\hG$; the leak law of $\GS$ with constant tending to $1$ in slip and energy (test P at three penalties, and $D$ in test sph, with slip $0.98$ and leak $0.93$ at $N=4$); and the removal of the leak by $\CNS$ (the leak correlation of $\CNS$ in test sph is $-0.38$ at $N=2$ and $-0.06$ at $N=4$, against $0.60$ and $0.87$ for $\GS$).

\emph{Not confirmed: the rates.} Only two sphere levels fit in memory, and they are pre-asymptotic. The observed $\hG$-rates of $2.4$--$2.5$ for both methods in tests sph and rot are dominated by the $P_3$ approximation error (the interpolation error drops from $1.63$ to $0.19$, and at $N=2$ it is larger than the finite element error) and by the $O(\hG^2)$ trace mismatch of $\ut$ on $\Gh$. At $N=4$ the finite element error is already $1.7$--$2.2$ times the interpolation error, so the geometric $\hG^{3/2}$ part is taking over; measuring its rate needs $N\ge8$. In test sph the $\GS$ leak ($D$, $0.064$ in $H^1$ at $N=4$) is still small compared with the geometric error ($0.33$), so the rate $1$ of $\GS$ and the rate $3/2$ of $\CNS$ in total $H^1$ error will separate only at much finer levels. The drift of $h/\hG$ also affects rates in $\hG$. We therefore make no claim about asymptotic rates on the sphere from these computations; the commands for $N=6,\dots,12$ (up to $1.85$ million unknowns) are in Appendix~\ref{app:repro}.
